\question[10]
On the figure below, identify the eight terrestrial biogeographic regions. Draw horizontal lines that show the approximate north/south boundaries between each of the regions. You do not need to draw east/west boundaries. Identify the location of Wallace's Line.

\ifprintanswers
	\includegraphics[width=\textwidth]{globalmap_answer}	
\else
	\includegraphics[width=\textwidth]{GlobalMap}
\fi



\question
\begin{parts}
\part[2] 
Below is the distribution of five families of freshwater crabs. To the left of the image, write the \emph{most likely} type of distribution displayed by the crabs.\vspace{2ex}

\begin{minipage}{0.35\textwidth}%
\ifprintanswers
	\textbf{Gondwanan}
\else
	\hspace{0.35\textwidth}
\fi
\end{minipage}\begin{minipage}{0.45\textwidth}%
	\hfill\includegraphics{crabdistrib}
\end{minipage}

\vspace{2ex}

\part[8] 
Klaus et al. (2011) examined whether the evolutionary relationships of the freshwater crabs were consistent with the distribution you just named.  Below left is a dendrogram that shows the order of the breakup of the land masses.  Below right is a phylogeny of the relationships among four of the crab families. Comparing the phylogeny to the dendrogram, explain whether the relationship of the crabs is consistent with a vicariant distribution. Provide two specific examples from the two trees to justify your answer. (Numbers in trees indicate time of divergence in millions of years ago (MYA). \textbf{NOTE TO SELF: Include earliest known fossil forms for taxa, wehre known?}\vspace{2ex}

\includegraphics{crabgondwana.png}
\end{parts}

	\begin{solution}
	
	Overall, no. Consistent if you allow for Dispersal from Indomalaysia to Eurasia, placing most Indomalaysian taxa as basal to clade of Neotropics and Africa.  However, still have problems with Socotra and Madagascar.  Madagascar taxa are more closely related to African taxa. Madagascar taxa should be more closely related to Indomalaysian taxa. Madagascar taxa diverged from African taxa circa 76–73 MYA, well after the separation of Madagascar (and India) from Africa. Neotropical taxa more closely related to Socotra and Eurasia taxa rather than African taxa. Seychelles more closely related to one clade of African taxa rather than Indomalaysian/Madagascar taxa.
	\end{solution}



\question
\begin{parts}
\part[5]
The figure below shows the similarity of coastal bivalves (clams and mussels) between North America and Europe. Briefly explain how and why bivalve similarity between North America and Europe is changing over time. (MYA = millions of years ago.)\vspace{2ex}

\begin{minipage}{0.47\textwidth}%
\ifprintanswers
	\textbf{Similarity decreases with distance between the continents because continued speciation over time causes the continental faunas to become increasingly different. }
\else
	\hspace{0.47\textwidth}
\fi
\end{minipage}\begin{minipage}{0.45\textwidth}%
	\hfill\includegraphics{atlanticbasin}
\end{minipage}

\vspace{2ex}

\part[10]
Briefly define or explain seafloor spreading and ridge push, and then describe how these processes together contribute to the pattern displayed in the figure above. You may draw a picture to aid your explanation.
	\begin{solution}
	Ridge push, seafloor spreading is moving the basins apart.
	\end{solution}

\end{parts}
